% Build: latexmk -pdf main.tex
\documentclass[11pt,a4paper]{article}
\usepackage[T1]{fontenc}
\usepackage[utf8]{inputenc}
\usepackage{libertine}
\usepackage[libertine]{newtxmath}
\usepackage[a4paper,margin=2.5cm]{geometry}
\usepackage{array}
\usepackage{microtype,graphicx,booktabs,enumitem,xcolor,titlesec,caption}
\usepackage[numbers,sort&compress]{natbib}
\usepackage[hidelinks]{hyperref}
\usepackage{xurl}
\urlstyle{same}

\definecolor{linkblue}{HTML}{1F5FBF}
\hypersetup{colorlinks=true,citecolor=linkblue,linkcolor=linkblue,urlcolor=linkblue}
\setlength{\parindent}{0pt}
\setlength{\parskip}{0.55em}
\linespread{1.05}
\titleformat{\section}{\large\sffamily\bfseries}{\thesection}{0.7em}{}
\titlespacing*{\section}{0pt}{1.3em}{0.4em}
\titleformat{\subsection}{\normalsize\sffamily\bfseries}{}{0pt}{}
\titlespacing*{\subsection}{0pt}{0.9em}{0.2em}
\captionsetup{font=small,labelfont=bf,labelsep=period}
\setlist[description]{leftmargin=0pt,itemsep=0.2em,topsep=0.2em,font=\normalfont\bfseries}
\setlength{\bibsep}{0.25em}
\renewcommand{\bibfont}{\small\raggedright}
\makeatletter
\setlength{\@fptop}{0pt}\setlength{\@fpsep}{1.8em}\setlength{\@fpbot}{0pt plus 1fil}
\makeatother

\setlength{\textfloatsep}{10pt plus 2pt minus 2pt}\setlength{\floatsep}{8pt plus 2pt minus 2pt}\setlength{\intextsep}{8pt plus 2pt minus 2pt}   % tighter space around floats

% Marker for numbers or figures that still have to be measured and filled in.
\newcommand{\todo}[1]{\textcolor{red}{\textbf{[TODO: #1]}}}

\begin{document}
\emergencystretch=3em

{\large\scshape DD2306 GPU Programming, Homework 2}\par\vspace{0.6em}
{\raggedright\Large\sffamily\bfseries Report for HW2, Group N\par}
\vspace{0.8em}
{\large Changhe You}\quad{\small\href{mailto:changhe@kth.se}{changhe@kth.se}}\par
\vspace{0.2em}

\textbf{Environment.} NVIDIA H100 80GB HBM3 (compute capability 9.0), compiled with
\texttt{nvcc -O3 -arch=sm\_90 -lineinfo}. All times are wall-clock CPU timer readings
from \texttt{gettimeofday} unless stated otherwise.

\section{Question 1: 1D Convolution}

\subsection{Implementation}

The mask has fixed width $W=5$, and out-of-boundary input elements are treated as $0$.
For an output index $i$,
\[
  P[i]=\sum_{j=0}^{W-1} N[i+j-\lfloor W/2\rfloor]\,M[j].
\]
The \emph{basic} kernel assigns one thread to each output element and reads
$N$ and $M$ directly from global memory, skipping indices outside $[0,n)$.
The \emph{tiled} kernel assigns one block to $T=\texttt{TILE\_WIDTH}$ consecutive outputs.
The block first cooperatively loads $T+W-1$ inputs (the tile plus a halo of
$\lfloor W/2\rfloor=2$ elements on each side) into shared memory, writes $0$ for
halo elements outside the array, calls \texttt{\_\_syncthreads()}, and then every
thread computes its output from shared memory. Both kernels sum the five terms in
the same order. Results of the two kernels are compared on the host with the
tolerance $|a-b|\le 10^{-5}$, because floating-point results may differ in the last
bits (different rounding or fused multiply-add choices), so exact equality is not a
reliable test.

\subsection{Name 3 application domains where convolutions are used}

\begin{description}
  \item[Image and video processing.] Blurring, sharpening and edge detection apply a small
    kernel over the pixels (the 2D version of this exercise).
  \item[Digital signal processing.] Audio filtering, smoothing of sensor data and
    finite-impulse-response (FIR) filters are 1D convolutions.
  \item[Deep learning.] Convolutional neural networks use convolutions as their core
    layer for vision, speech and time-series models.
\end{description}

\subsection{Global memory reads and writes}

Let $n$ be the array length, $W=5$ the mask width, and $T$ the tile width.
Counts refer to accesses in the code, not to DRAM traffic; in practice the five
elements of $M$ are served by the caches.

\begin{description}
  \item[Basic kernel.] Each thread reads, for every valid $j$, one element of $N$ and
    one element of $M$. Only $6$ terms are lost at the two array ends
    ($2+1$ on the left, $1+2$ on the right), so $N$ is read $5n-6$ times and $M$ is
    read $5n-6$ times: $10n-12\approx 10n$ reads in total. Each thread writes one
    output, giving $n$ writes.
  \item[Tiled kernel (our implementation).] Each block loads $T+4$ input elements, and
    there are $\lceil n/T\rceil$ blocks, so $N$ is read about $n+4n/T$ times (slightly
    fewer at the array ends where ghost cells are not loaded). Our kernel still reads
    $M[j]$ from global memory in every thread, which gives $5n$ reads of $M$. The total
    is about $6n+4n/T$ reads and $n$ writes.
\end{description}

Tiling therefore reduces the reads of the input $N$ from about $5n$ to about
$n(1+4/T)$, i.e.\ close to one read per element. Because $M$ is not cached in shared or
constant memory in our kernel, the total only drops from $10n$ to about $6n$. Keeping
$M$ in constant memory or in shared memory would bring the tiled total towards
$n(1+4/T)$ reads. For example, with $n=16384$ and $T=128$ the counts are
$163{,}828$ reads for the basic kernel and $98{,}812$ reads for the tiled kernel
($16{,}892$ of them for $N$), both with $16{,}384$ writes.

\subsection{Results for $N=1024,\dots,16384$}

Figure~\ref{fig:q1-output} shows the program output and Table~\ref{tab:q1} lists the
kernel times. All five runs print \texttt{PASS}, and the first outputs match a hand
calculation, for example
$P[0]=0.840188\cdot1.0+0.394383\cdot0.5+0.783099\cdot0.25=1.233154$.

\begin{figure}[t]
  \centering
  \IfFileExists{fig/q1_output.png}{%
    \includegraphics[width=\linewidth]{fig/q1_output.png}}{%
    \fbox{\parbox{0.9\linewidth}{\centering\todo{insert screenshot fig/q1\_output.png of the runs for $N=1024$ to $16384$}}}}
  \caption{Program output for $N=1024,2048,4096,8192,16384$.}
  \label{fig:q1-output}
\end{figure}

\begin{table}[t]
  \centering
  \caption{CPU-timer kernel times with $\texttt{TILE\_WIDTH}=128$ and block size $256$ for the basic kernel.}
  \label{tab:q1}
  \begin{tabular}{rrr}
    \toprule
    $N$ & basic ($\mu$s) & tiled ($\mu$s) \\
    \midrule
    1024  & 94  & 17 \\
    2048  & 94  & 16 \\
    4096  & 97  & 18 \\
    8192  & 93  & 17 \\
    16384 & 103 & 18 \\
    \bottomrule
  \end{tabular}
\end{table}

These numbers should not be read as ``tiled is five times faster''. The times are almost
independent of $N$, so they measure fixed launch overhead and not computation: even
$N=16384$ is only 64 tiled blocks, fewer than the 132 SMs of an H100. In addition, the
basic kernel is launched first and carries the one-time first-launch cost
(the first CUDA call also takes about $0.2$--$0.3$\,s for context creation, visible in
the ``Allocated device memory'' line). A fair comparison needs a warm-up launch,
\texttt{cudaEvent} timing, repeated runs, and larger $N$.
\todo{redo timing with warm-up, cudaEvent and repeats; replace Table~\ref{tab:q1}}

\subsection{Choosing the tile size}

\todo{run TILE\_WIDTH $=32,64,128,256,512,1024$ with the improved timing and report the measurements}

The considerations for the choice are:
\begin{itemize}[leftmargin=1.5em,itemsep=0.1em]
  \item \textbf{Halo overhead.} A block loads $T+4$ elements for $T$ outputs, so the
    redundant fraction is $4/T$: $12.5\%$ for $T=32$, but only $3\%$ for $T=128$.
    Larger tiles amortise the halo better.
  \item \textbf{Block size.} The block has $T$ threads, so $T$ should be a multiple of the
    warp size $32$ and at most $1024$.
  \item \textbf{Shared memory and occupancy.} A tile uses $4(T+4)$ bytes, which is tiny
    (about 0.5\,KB for $T=128$), so shared memory does not limit occupancy here. The block
    size does: very large $T$ allows few resident blocks per SM and, for small $N$,
    leaves SMs idle because there are too few blocks.
  \item \textbf{Parallelism for small $N$.} The number of blocks is $\lceil N/T\rceil$;
    with $N=1024$ and $T=1024$ there is a single block on a single SM.
\end{itemize}
\todo{state the chosen optimal tile size and justify it with the measured data}

\subsection{Profiling the tiled implementation}

\todo{Nsight Systems: timeline and kernel/memcpy breakdown}\par
\todo{Nsight Compute: performance hotspot (memory vs.\ compute), shared memory utilisation, achieved occupancy}\par
\todo{effect of changing the shared memory carveout (\texttt{cudaFuncSetAttribute} with \texttt{cudaFuncAttributePreferredSharedMemoryCarveout}) on runtime}

\section{Question 2: Tiled Matrix Multiplication}
\todo{not done yet}

\section{Question 3: Accelerating Llama2 Inference}
\todo{not done yet}

\end{document}
